\documentclass[10pt,a4paper]{amsart}
\usepackage[margin=19mm]{geometry}
\usepackage{amsmath,amssymb,mathtools}
\usepackage[scr=rsfs,cal=euler]{mathalfa}
\usepackage{xfrac}
\usepackage[unicode,hidelinks,bookmarks=false]{hyperref}
\newcommand{\ie}{i.e.\ }
\newcommand{\cf}{cf.\ }
\newcommand{\resp}{respectively\ }
\newcommand{\Subsection}{\subsection}
\providecommand{\nobreakdash}{\nobreak\hbox{-}}
\setlength{\emergencystretch}{2em}
\multlinegap0pt
\usepackage{bm}
\usepackage{bbm}
\usepackage{dsfont}
\usepackage{xspace}
\usepackage{cancel}
\usepackage{mathtools}
\usepackage{amsthm}
\makeatletter
\@ifundefined{theorem}{\newtheorem{theorem}{Theorem}}{}
\@ifundefined{proposition}{\newtheorem{proposition}[theorem]{Proposition}}{}
\makeatother
\newcommand{\Z}{{\mathbb{Z}}}
\title[The numerical range of periodic banded Toeplitz operators]{The numerical range of periodic banded Toeplitz operators}
\author{Benjam\'in A. Itz\'a-Ortiz, Rub\'en A. Mart\'inez-Avenda\~no, Hiroshi Nakazato}
\date{Source: arXiv:2308.12353v1 (CC BY 4.0)}
\begin{document}
\maketitle

\noindent\emph{Source and licence.} The numerical range of periodic banded Toeplitz operators. Version arXiv:2308.12353v1 (CC BY 4.0), distributed under the Creative Commons Attribution 4.0 International licence, as recorded in the arXiv licence metadata for this exact version and shown on the abstract page https://arxiv.org/abs/2308.12353v1. Full rights notice: licenses/arxiv-2308.12353-CC-BY-4.0.txt.\par\smallskip

\noindent\textit{Editorial scope.} Counted unit: the Proposition whose source label is \texttt{prop:Cmu} in the article (source0.tex lines 331--336, source label \texttt{prop:Cmu}), delivered together with the article's own complete original proof of it (source0.tex lines 337--428, one \texttt{proof} environment of 92 lines). No mathematics is written, repaired or re-derived: statement and proof are the article's own text line for line. Required context retained uncounted: none. Every definition, display and label of those lines is retained unchanged. Omitted from this package and not redistributed: 1--330, 429--833. Nothing inside a retained excerpt is omitted or altered: every retained body is delivered exactly as the article's lines are written, so no omission receipt of any kind accompanies this package. Citations: the retained excerpts carry no citation command at all, so no bibliography entry of the article is transcribed and no reference is redistributed. Reference adaptation: none was needed and none was made; every reference inside the delivered unit resolves inside it, so no unresolved reference or citation is produced. Printed slips kept literal: no printed word, symbol, hypothesis or number of the counted statement or of its proof was altered. Layout and class: the article's own class is \texttt{amsart} with packages including amssymb, amsbsy, amscd, enumerate, setspace, tikz, bm, cite, graphicx, hyperref, nicematrix, colortbl, mathtools, kpfonts; the wrapper is the lane's pinned \texttt{amsart} editorial wrapper at 10pt on A4, which restates the theorem-like environments the retained text needs and adds the article's own macro definitions that the retained text needs (\texttt{\textbackslash Z}), with extra wrapper packages (bm, bbm, dsfont, xspace, cancel, mathtools). The delivered numbering is the unit's own. Assets: no retained span contains a figure environment, a graphics-inclusion command, a PDF-page inclusion, a table or a TikZ picture; no image file is copied into this package, the package declares no asset dependency and no image file is redistributed with it. Licence: version arXiv:2308.12353v1 of this article is distributed under the Creative Commons Attribution 4.0 International licence, as recorded in the arXiv licence metadata for this exact version and shown on the abstract page https://arxiv.org/abs/2308.12353v1. A full rights notice is delivered at licenses/arxiv-2308.12353-CC-BY-4.0.txt. Characters: every retained excerpt is pure ASCII, so the delivered engine, XeLaTeX with its default Latin Modern text font, reports no missing character.
\begin{proposition}\label{prop:Cmu}
Let $n$, $s$ and $\mu$ be as above. Then $C_\mu$ is unitarily equivalent to a block diagonal matrix. More precisely,   
\[
  U^{*}C_{\mu} U=\Phi(0) \oplus \Phi(2\pi/s) \oplus \Phi(4\pi/s) \oplus \cdots \oplus \Phi(2(s-1)\pi/s).
  \]
\end{proposition}

\begin{proof}
Let $p_1, p_2, q_1, q_2$ be integers such that $0 \leq q_1, q_2 \leq s-1$ and $0\leq p_1, p_2 \leq n$. Observe that
\begin{align*}
   \langle C_{\mu} f_{p_1, q_1}, f_{p_2, q_2} \rangle 
   & =\frac{1}{s} \sum_{u=0}^{s-1} \sum_{v=0}^{s-1} \langle C_{\mu} \rho^{u q_1} e_{p_1+u (n+1)}, \rho^{v q_2} e_{p_2+v (n+1)} \rangle \\
   &=\frac{1}{s}\sum_{u=0}^{s-1} \sum_{v=0}^{s-1}  \rho^{u q_1} \rho^{-v q_2}\langle C_{\mu} e_{p_1+u (n+1)}, e_{p_2+v (n+1)} \rangle \\ 
   &=\frac{1}{s}\sum_{u=0}^{s-1} \sum_{v=0}^{s-1}  \rho^{u q_1} \rho^{-v q_2}  \sum_{w \in {\mathbb Z}}  \langle 
      T e_{p_1 +u (n+1) +w \mu}, e_{p_2 +v (n+1)} \rangle \\
&=\frac{1}{s}\sum_{u=0}^{s-1} \sum_{v=0}^{s-1} \rho^{u q_1} \rho^{-v q_2} \sum_{w \in {\mathbb Z}}  \langle
      T e_{p_1 +(u-v) (n+1) +w \mu}, e_{p_2} \rangle,   
\end{align*}
where the last equality follows from the $(n+1)$-periodicity of $T$.

Making the change of variables $u=v+r$, we obtain
\begin{align*}
   \langle C_{\mu} f_{p_1, q_1}, f_{p_2, q_2} \rangle 
   & =\frac{1}{s}\sum_{v=0}^{s-1} \sum_{r=-v}^{-v+s-1}  \rho^{(v+r)q_1} \rho^{-v q_2}  \sum_{w \in {\mathbb Z}}  \langle
      T e_{p_1 + r (n+1) +w \mu}, e_{p_2} \rangle.
    \end{align*}
    
For each $v=0, 1, 2, \dots, s-1$, we consider the inner sum of the above expression and rewrite it as
\begin{align*}
    \sum_{r=-v}^{-v+s-1}  & \rho^{(v+r)q_1} \rho^{-v q_2}  \sum_{w \in {\mathbb Z}}   \langle
      T e_{p_1 + r (n+1) +w \mu}, e_{p_2} \rangle \\
      &= \sum_{r=-v}^{-1}  \rho^{(v+r)q_1} \rho^{-v q_2}  \sum_{w \in {\mathbb Z}}  \langle
      T e_{p_1 + r (n+1) +w \mu}, e_{p_2} \rangle + \sum_{r=0}^{-v+s-1}  \rho^{(v+r)q_1} \rho^{-v q_2}  \sum_{w \in {\mathbb Z}}  \langle
      T e_{p_1 + r (n+1) +w \mu}, e_{p_2} \rangle.
\end{align*}    
Making another change of variables in the first summand above, we obtain
\begin{align*}
    \sum_{r=-v}^{-1} & \rho^{(v+r)q_1} \rho^{-v q_2}  \sum_{w \in {\mathbb Z}}  \langle
      T e_{p_1 + r (n+1) +w \mu}, e_{p_2} \rangle + \sum_{r=0}^{-v+s-1}  \rho^{(v+r)q_1} \rho^{-v q_2}  \sum_{w \in {\mathbb Z}}  \langle
      T e_{p_1 + r (n+1) +w \mu}, e_{p_2} \rangle \\
      &= \sum_{r=s-v}^{s-1}  \rho^{(v+r-s)q_1} \rho^{-v q_2}  \sum_{w \in {\mathbb Z}}  \langle
      T e_{p_1 + (r-s) (n+1) +w \mu}, e_{p_2} \rangle + \sum_{r=0}^{-v+s-1}  \rho^{(v+r)q_1} \rho^{-v q_2}  \sum_{w \in {\mathbb Z}}  \langle
      T e_{p_1 + r (n+1) +w \mu}, e_{p_2} \rangle \\
&= \sum_{r=s-v}^{s-1}  \rho^{(v+r)q_1} \rho^{-v q_2}  \sum_{w \in {\mathbb Z}}  \langle
      T e_{p_1 + r(n+1) +(w-1) \mu}, e_{p_2} \rangle + \sum_{r=0}^{-v+s-1}  \rho^{(v+r)q_1} \rho^{-v q_2}  \sum_{w \in {\mathbb Z}}  \langle
      T e_{p_1 + r (n+1) +w \mu}, e_{p_2} \rangle \\
      &= \sum_{r=s-v}^{s-1}  \rho^{(v+r)q_1} \rho^{-v q_2}  \sum_{w \in {\mathbb Z}}  \langle
      T e_{p_1 + r(n+1) +w \mu}, e_{p_2} \rangle + \sum_{r=0}^{-v+s-1}  \rho^{(v+r)q_1} \rho^{-v q_2}  \sum_{w \in {\mathbb Z}}  \langle
      T e_{p_1 + r (n+1) +w \mu}, e_{p_2} \rangle \\
      &= \sum_{r=0}^{s-1}  \rho^{(v+r)q_1} \rho^{-v q_2}  \sum_{w \in {\mathbb Z}}  \langle
      T e_{p_1 + r(n+1) +(w+1) \mu}, e_{p_2} \rangle.
\end{align*}
    
Hence
\begin{equation}\label{eq:Cmfpq}
\begin{split}
   \langle C_{\mu} f_{p_1, q_1}, f_{p_2, q_2} \rangle 
   & = \frac{1}{s} \sum_{v=0}^{s-1} \sum_{r=0}^{s-1}  \rho^{(v+r)q_1} \rho^{-v q_2}  \sum_{w \in {\mathbb Z}}  \langle
      T e_{p_1 + r(n+1) +(w+1) \mu}, e_{p_2} \rangle \\
      &= \frac{1}{s}\sum_{w \in {\mathbb Z}}
       \sum_{r=0}^{s-1} 
       \sum_{v=0}^{s-1} \rho^{(v+r)q_1} \rho^{-v q_2}  \langle
      T e_{p_1 + r(n+1) +(w+1) \mu}, e_{p_2} \rangle  \\
      &=\frac{1}{s}\sum_{w \in {\mathbb Z}}
       \sum_{r=0}^{s-1} \rho^{r q_1}  \langle
      T e_{p_1 + r(n+1) +(w+1) \mu}, e_{p_2} \rangle 
       \sum_{v=0}^{s-1} \rho^{v(q_1-q_2)}.
       \end{split}
   \end{equation}

Now, assume $q_1 \neq q_2$. Since this implies that  $\sum_{v=0}^{s-1} \rho^{v(q_1-q_2)}=0$, equation \eqref{eq:Cmfpq} gives
\[
\langle C_{\mu} f_{p_1, q_1}, f_{p_2, q_2} \rangle =0.
\]

If $q_1=q_2$, then $\sum_{v=0}^{s-1} \rho^{v(q_1-q_2)}=s$ and hence equation \eqref{eq:Cmfpq} gives
\begin{align*}
\langle C_{\mu} f_{p_1, q_1}, f_{p_2, q_2} \rangle 
&=\sum_{w \in {\mathbb Z}}
       \sum_{r=0}^{s-1} \rho^{r q_1}  \langle
      T e_{p_1 + r(n+1) +(w+1) \mu}, e_{p_2} \rangle \\
&=\sum_{w \in {\mathbb Z}}
       \sum_{r=0}^{s-1} \rho^{r q_1}  \langle
      T e_{p_1 + (r+(w+1)s)(n+1)}, e_{p_2} \rangle \\
&=\sum_{w \in {\mathbb Z}}
       \sum_{r=0}^{s-1} \rho^{(r+(w+1)s)q_1}  \langle
      T e_{p_1 + (r+(w+1)s)(n+1)}, e_{p_2} \rangle,
\end{align*}
since $\rho^s=1$.

Observe that when $r$ ranges from $0$ to $s-1$ and $w$ ranges in $\Z$, the number $r+(w+1)s$ ranges over all the integers, and hence the previous expression can be written as
\[
\langle C_{\mu} f_{p_1, q_1}, f_{p_2, q_2} \rangle = \sum_{u \in {\mathbb Z}} \rho^{u q_1}  \langle
      T e_{p_1 + u (n+1)}, e_{p_2} \rangle,
\]
which equals the $(p_2,p_1)$ entry of the matrix $\Phi\left(\frac{2\pi q_1}{s}\right)$.

The above computations show that the matrix of $C_\mu$ with respect to the orthogonal basis $\{ f_{p,q} \, : \, 0 \leq p \leq n, \ 0  \leq q \leq s-1\}$ is block-diagonal and each of the $s$ blocks on the diagonal is of the form $\Phi\left(\frac{2\pi q}{s}\right)$, for $q=0, 1, \dots, s-1$. This proves the desired result.
\end{proof}

\end{document}
